\documentclass[a4paper]{scrartcl}

% --- Packages ---
\usepackage[a4paper, left=1.8cm, right=1.8cm, top=2.5cm, bottom=2cm]{geometry}
\usepackage[utf8]{inputenc}
\usepackage[T1]{fontenc}
\usepackage{amsmath, amsfonts, amssymb}
\usepackage{graphicx}
\usepackage{svg}
\usepackage{float}
\usepackage{longtable}
\usepackage{booktabs}
\usepackage{adjustbox}
\usepackage{makecell}
\usepackage{threeparttable}
\usepackage{array}
\usepackage{siunitx}
\usepackage{multicol}
\usepackage{wrapfig}
\usepackage{caption}
\usepackage{algorithm}
\usepackage{algorithmic}
\usepackage{scrlayer-scrpage}
\usepackage{xcolor}
\usepackage{changepage}

\renewcommand{\arraystretch}{1.4}
\pagestyle{scrheadings}
\automark{section}
\ohead{\pagemark}
\cfoot{}

% --- Document Info ---
\title{SCF-PDP Programming Exercise Report: Assignment 2}
\subtitle{Heuristic Optimization Techniques, WS 2025}
\author{
    Daniel Levin (12433760)
}
\date{}

% --- Custom Environments ---
\newtheorem{ex}{Exercise}

\begin{document}
\maketitle

% ============================================================
% 1. Introduction / Overview
% ============================================================
\section{Introduction}
This report presents our work on Assignment 2 of the SCF-PDP programming exercise. Building upon the algorithms developed in Assignment 1 (construction heuristics, local search, VND, GRASP, and Simulated Annealing), we implement two advanced metaheuristics, investigate alternative fairness measures, perform automated parameter tuning, and conduct statistical comparisons.

% ============================================================
% 2. General Setup (reuse/extend from Assignment 1)
% ============================================================
\section{General Setup}

\subsection{Implementation Environment}
% TODO: Update if any changes from Assignment 1
The implementation continues in Python 3.13 with the following libraries:
\begin{itemize}
  \item \texttt{numpy} -- efficient array operations and numerical computations
  \item \texttt{scipy} -- distance matrix computation and statistical functions
  \item \texttt{pymhlib} -- metaheuristic framework
  \item \texttt{pandas} -- result aggregation
  \item \texttt{matplotlib} -- visualization
  % TODO: Add any new libraries used
\end{itemize}

\subsection{Extensions from Assignment 1}
% TODO: Describe any changes to core classes (SCFPDPInstance, SCFPDPSolution, Route)
% TODO: Describe new data structures or utilities added

% ============================================================
% 3. Advanced Metaheuristic 1
% ============================================================
\section{Advanced Metaheuristic 1: [Algorithm Name]}
% Choose from: Evolutionary/Genetic Algorithm, Ant Colony Optimization,
% (Adaptive) Large Neighborhood Search, or Hybrid Algorithm

\subsection{Idea}
% TODO: Describe the main concept and motivation for choosing this algorithm

\subsection{Algorithm Design}
% TODO: Describe the algorithm structure
% Include: representation, operators, key mechanisms

\subsection{Implementation Details}
% TODO: Describe specific implementation choices
% Include: data structures, key functions

\subsection{Parameters}
% TODO: List all parameters with initial/default values
% Example table:
% \begin{table}[H]
% \centering
% \caption{Algorithm 1 Parameters}
% \begin{tabular}{l l l}
% \toprule
% Parameter & Description & Default Value \\
% \midrule
% ... & ... & ... \\
% \bottomrule
% \end{tabular}
% \end{table}

% ============================================================
% 4. Advanced Metaheuristic 2
% ============================================================
\section{Advanced Metaheuristic 2: [Algorithm Name]}
% Choose a different algorithm from the list above

\subsection{Idea}
% TODO: Describe the main concept and motivation for choosing this algorithm

\subsection{Algorithm Design}
% TODO: Describe the algorithm structure

\subsection{Implementation Details}
% TODO: Describe specific implementation choices

\subsection{Parameters}
% TODO: List all parameters with initial/default values

% ============================================================
% 5. Fairness Measure Investigation (Task 2)
% ============================================================
\section{Fairness Measure Investigation}

\subsection{Alternative Fairness Measures}
We investigate two alternative fairness measures in addition to the Jain fairness index:

\subsubsection{Max-Min Fairness}
The max-min fairness measure is defined as:
\[
M = \frac{\min_{k \in K} d(R_k)}{\max_{k \in K} d(R_k)}
\]
% TODO: Describe interpretation and properties

\subsubsection{Gini Coefficient (Inverted)}
The inverted Gini coefficient is defined as:
\[
G = 1 - \frac{\sum_{k' \in K} \sum_{k'' \in K} |d(R_{k'}) - d(R_{k''})| }{2 n_K \sum_{k' \in K} d(R_{k'})}
\]
% TODO: Describe interpretation and properties

\subsection{Implementation}
% TODO: Describe how the fairness measures were implemented

\subsection{Experimental Results}
% TODO: Compare how objective values change with different fairness measures
% TODO: Analyze how the number of stops per route changes
% Include tables and figures

\subsection{Analysis}
% TODO: Discuss findings and insights about fairness measure impact

% ============================================================
% 6. Automated Parameter Tuning (Task 3)
% ============================================================
\section{Automated Parameter Tuning}

\subsection{Tuning Methodology}
% TODO: Describe the parameter tuning approach used
% Options: Grid search, Random search, irace, ParamILS, SMAC, etc.

\subsection{Tuning Setup}
% TODO: Describe:
% - Parameter ranges explored
% - Tuning instances used
% - Evaluation metrics
% - Computational resources (e.g., cluster usage)

\subsection{Algorithm 1 Tuning Results}
% TODO: Present tuning results for first algorithm
% Include: parameter configurations tested, best found parameters
% Show results per instance size

\subsection{Algorithm 2 Tuning Results}
% TODO: Present tuning results for second algorithm

\subsection{Instance Size Scalability}
% TODO: Report which instance sizes can be solved in reasonable time

% ============================================================
% 7. Experimental Results on Test Instances (Task 4)
% ============================================================
\section{Experimental Results}
ABCD

\subsection{Experimental Setup}
% TODO: Describe test instances, number of runs, time limits

\subsection{Best Algorithm from Assignment 1}
% TODO: Present results of the best algorithm from Assignment 1 on test instances
% Include: mean objective, std dev, runtime

\subsection{Best Algorithm from Assignment 2}
% TODO: Present results of the better tuned algorithm from this assignment
% Justify the choice if results were inconclusive

\subsection{Comparison Across Instance Sizes}
% TODO: Show how algorithms scale with instance size
% Include tables and figures

\subsection{Solution Quality Over Time}
% TODO: Include convergence plots showing quality over iterations/time
% \begin{figure}[H]
%     \centering
%     \includegraphics[width=1\linewidth]{convergence_plot.png}
%     \caption{Solution Quality Over Time}
%     \label{fig:convergence}
% \end{figure}

% ============================================================
% 8. Statistical Comparison (Task 5)
% ============================================================
\section{Statistical Comparison}

\subsection{Methodology}
% TODO: Describe statistical tests used (e.g., Wilcoxon signed-rank test)
% Reference the Statistical Testing Jupyter notebook

\subsection{Comparison: Best from Assignment 1 vs. Best from Assignment 2}
% TODO: Present statistical test results
% Include: p-values, effect sizes, conclusions

\subsection{Variance Analysis}
% TODO: Include box plots or violin plots showing solution variance
% \begin{figure}[H]
%     \centering
%     \includegraphics[width=1\linewidth]{boxplot_comparison.png}
%     \caption{Solution Quality Distribution Comparison}
%     \label{fig:boxplot}
% \end{figure}

\subsection{Summary of Statistical Findings}
% TODO: Summarize which algorithm is significantly better (if any)

% ============================================================
% 9. Conclusion
% ============================================================
\section{Conclusion}
% TODO: Summarize main findings and contributions

\subsection{Key Learnings}
% TODO: Reflect on experiences from both programming assignments
% What worked well? What was challenging?

\subsection{Algorithm Design Insights}
% TODO: Discuss insights about metaheuristic design for SCF-PDP

\subsection{Future Improvements}
% TODO: What would you do differently or explore further?

% ============================================================
% Appendices (optional)
% ============================================================
% \appendix
% \section{Additional Plots}
% \section{Parameter Tables}
% \section{Code Snippets}

% ============================================================
% Bibliography (optional)
% ============================================================
% \bibliographystyle{plain}
% \begin{thebibliography}{9}
% \bibitem{}
% \textit{},
% \url{}
% \end{thebibliography}

\end{document}